\chapter*{Заключение}
\addcontentsline{toc}{chapter}{Заключение}
\label{ch:conclusion}

В ходе выполнения лабораторной работы~\cfgLabNumber{} получены
практические навыки разработки на JavaScript, охватывающие путь от
базовых конструкций языка до публикации полноценного веб-приложения.

Решены задачи всех четырёх уроков методички:

\begin{itemize}
  \item на уроке~1 освоено объявление переменных
        (\texttt{const}, \texttt{let}), выбор типов данных
        (\texttt{string}, \texttt{number}) и вывод результатов в
        консоль через \texttt{console.log};
  \item на уроке~2 проработана работа с массивами и циклами:
        итерация \texttt{for} с условным выбором имени отправителя,
        фильтрация массива методом \texttt{filter} в связке с
        \texttt{String.prototype.includes};
  \item на уроке~3 реализован одностраничный сайт-каталог героев
        Marvel: данные получаются с публичного REST-API через
        \texttt{fetch}\,/\,\texttt{async-await}, разметка карточек и
        модальных окон формируется функциями
        \texttt{getCharacterCard[s]} и \texttt{getCharacterModal[s]} в
        связке с шаблонными строками; компоненты Bootstrap~5
        (\texttt{card}, \texttt{modal}, \texttt{spinner}) обеспечивают
        внешний вид и адаптивность без единой строки собственного CSS;
  \item на уроке~4 готовое приложение опубликовано на собственном
        домашнем сервере под обратным прокси \textbf{Caddy} с
        автоматическим выпуском TLS-сертификата Let's~Encrypt; адрес
        развёрнутого сайта --- \url{https://marvel.server34.netcraze.club}.
\end{itemize}

Дополнительные приёмы, не вынесенные в явные требования методички, но
применённые в проекте:
\begin{itemize}
  \item экранирование пользовательских данных через собственную
        функцию \texttt{escapeHtml} --- защита от XSS при
        использовании \texttt{innerHTML};
  \item принудительный переход URL обложек с \texttt{http} на
        \texttt{https} --- обход блокировки \textit{mixed content} на
        HTTPS-странице;
  \item обработка ошибки сети с информативным сообщением
        пользователю на месте спиннера;
  \item автоматизация развёртывания через идемпотентный bash-скрипт
        в репозитории \texttt{home-server}.
\end{itemize}

Полученные навыки --- основа любой клиентской разработки в вебе.
Концепции событийной модели DOM, асинхронных запросов через промисы
и шаблонных строк применимы как в проектах на чистом JavaScript, так
и в большинстве современных фреймворков (\textit{React}, \textit{Vue},
\textit{Svelte}), которые лишь дополняют изученный фундамент своими
абстракциями (виртуальный DOM, реактивность, компоненты).

\endinput
